% ============================================================================
%  CAPÍTULO III: DESARROLLO DEL TRABAJO
% ============================================================================

\section{Desarrollo del trabajo}

Este cap\'itulo documenta, paso a paso, la configuraci\'on del \textbf{DNS del
dominio} (servido por el Controlador de Dominio de Samba) y del \textbf{servidor
web} de la organizaci\'on (Nginx como \textit{proxy inverso}), sobre el
servidor \'unico \texttt{dc1}. Cada bloque describe desde la arquitectura de la
soluci\'on hasta las pruebas de resoluci\'on de nombres y de acceso al portal,
incluyendo la justificaci\'on t\'ecnica de cada una de las decisiones tomadas.

% ============================================================================
\section{Arquitectura general del par DNS+WEB}

El proyecto concentra todos los servicios en un \'unico servidor
(\texttt{192.168.0.2}), por lo que el DNS del dominio y la puerta web
comparten esa misma direcci\'on. La separaci\'on es \textbf{l\'ogica}, no
f\'isica: \citep{tanenbaum2012}

\begin{itemize}
  \item \textbf{DNS:} el Controlador de Dominio de \textbf{Samba 4} sirve la
        zona autoritativa \texttt{sudoers.lan} en el puerto 53 de forma
        \textbf{nativa} (fuera de Docker), usando el backend
        \textbf{SAMBA\_INTERNAL}.
  \item \textbf{WEB:} \textbf{Nginx} corre en un \textbf{contenedor Docker}
        (servicio \texttt{web} del \texttt{docker-compose.yml}) y publica los
        puertos 80/443 del host. Act\'ua como \textit{proxy inverso}: toda
        petici\'on \texttt{*.sudoers.lan} entra por \'el y es enrutada al
        servicio que corresponda \citep{nginxDocs}.
\end{itemize}

La tabla \ref{tab:mapa-nombres} resume los nombres gestionados por el DNS del
dominio y a qu\'e backend los dirige el servidor web.

\begin{table}[H]
  \centering
  \caption{Mapa de nombres gestionados por el DNS de \texttt{sudoers.lan}}
  \contenidoConFuente{%
    \begin{tablaAPA}
      \begin{tabular}{@{}p{4.6cm}p{6.2cm}c@{}}
        \toprule
        \textbf{Nombre} & \textbf{Backend del proxy} & \textbf{Registro} \\
        \midrule
        \texttt{sudoers.lan} / \texttt{web} / \texttt{www} & Portal est\'atico (Nginx) & A, CNAME \\
        \texttt{dc1.sudoers.lan} & DC (DNS, SMB) & A \\
        \texttt{print.sudoers.lan} & \texttt{tso-print:631} (CUPS) & A \\
        \texttt{portainer.sudoers.lan} & \texttt{tso-portainer:9000} & A \\
        \texttt{joaquin.sudoers.lan} & P\'agina del equipo (est\'atico) & A \\
        \texttt{david.sudoers.lan} & \texttt{192.168.0.51:8080} & A \\
        \texttt{nicolas.sudoers.lan} & \texttt{192.168.0.52:8080} & A \\
        \texttt{\_ldap.\_tcp.dc.\_msdcs} & \texttt{dc1}:389 (LDAP) & SRV \\
        \bottomrule
      \end{tabular}
    \end{tablaAPA}%
  }{Elaboraci\'on propia.}
  \label{tab:mapa-nombres}
\end{table}

% ============================================================================
\section{El DNS del dominio: Samba AD (SAMBA\_INTERNAL)}

\subsection{Autoridad de la zona y backend DNS}

Un dominio de Active Directory se apoya por completo en el DNS: los clientes
lo consultan para descubrir d\'onde est\'a cada servicio del dominio, tanto
los nombres comunes (\texttt{web}) como los registros de descubrimiento (SRV)
\citep{msft_activedirectory}. Por ello, durante el aprovisionamiento del DC se
eligi\'o expl\'icitamente que el propio Samba administre la zona:

\begin{lstlisting}[style=terminal]
$ samba-tool domain provision \
    --realm=SUDOERS.LAN --domain=SUDOERS \
    --server-role=dc --dns-backend=SAMBA_INTERNAL \
    --use-rfc2307 --host-name=dc1 --host-ip=192.168.0.2 \
    --option="dns forwarder = 192.168.0.1"
\end{lstlisting}

Con el backend \textbf{SAMBA\_INTERNAL} el DC es el servidor de nombres
autoritativo de la zona \texttt{sudoers.lan}: guarda los registros dentro del
directorio (no en archivos de zona BIND) y los actualiza din\'amicamente.
Adem\'as se le indic\'o un \textit{forwarder} hacia el router
(\texttt{192.168.0.1}) para resolver nombres externos de Internet
\citep{iscBIND9}. La zona queda registrada en \texttt{/etc/samba/smb.conf}
junto con el rol del servidor:

\begin{lstlisting}[style=terminal]
$ grep -E 'server role|dns|realm' /etc/samba/smb.conf
[global]
    server role = active directory domain controller
    realm = SUDOERS.LAN
    dns forwarder = 192.168.0.1
    dns update command = /usr/sbin/samba_dnsupdate
\end{lstlisting}

\capturaAPA{fig:smb-dns}{Zona \texttt{sudoers.lan} y rol de DC en
\texttt{smb.conf}}{captura-smb-dns.png}{Captura de pantalla mostrando la salida de \texttt{grep} sobre \texttt{/etc/samba/smb.conf} con el rol \texttt{active directory domain controller}, el \textit{realm} \texttt{SUDOERS.LAN} y la directiva \texttt{dns forwarder = 192.168.0.1}.}

\subsection{Registros A de los servicios}

Como todos los servicios corren en el \'unico servidor de la red, los
registros A de la zona apuntan, salvo el DC mismo, al \textbf{proxy inverso}
\texttt{192.168.0.2}. El listado completo de la zona se consulta con
\texttt{samba-tool dns}:

\begin{lstlisting}[style=terminal]
$ kinit administrator                    # ticket de administrador
$ samba-tool dns query dc1.sudoers.lan sudoers.lan @ ALL -k yes
Name=, Records=8, Children=0
 DC@SUDOERS.LAN: type=A, value=192.168.0.2
 web@SUDOERS.LAN: type=A, value=192.168.0.2
 joaquin@SUDOERS.LAN: type=A, value=192.168.0.2
 print@SUDOERS.LAN: type=A, value=192.168.0.2
 portainer@SUDOERS.LAN: type=A, value=192.168.0.2
 david@SUDOERS.LAN: type=A, value=192.168.0.2
 nicolas@SUDOERS.LAN: type=A, value=192.168.0.2
\end{lstlisting}

\capturaAPA{fig:zona-dns}{Zona DNS completa de \texttt{sudoers.lan} consultada
con \texttt{samba-tool dns}}{captura-zona-dns.png}{Captura de pantalla mostrando la salida de \texttt{samba-tool dns query dc1.sudoers.lan sudoers.lan @ ALL} con los registros A de \texttt{dc1}, \texttt{web}, \texttt{print}, \texttt{portainer} y de los administradores, todos hacia \texttt{192.168.0.2}.}

\subsection{Registro A del Controlador de Dominio}

El registro \texttt{A} es la tabla nombre$\rightarrow$IP del dominio: toda
consulta al DNS de la red empieza pidiendo \texttt{dc1.sudoers.lan}. La
resoluci\'on directa se verifica con \texttt{dig}:

\begin{lstlisting}[style=terminal]
$ dig web.sudoers.lan
;; ANSWER SECTION:
web.sudoers.lan. 900 IN A 192.168.0.2
\end{lstlisting}

\capturaAPA{fig:dig-web}{Resoluci\'on de \texttt{web.sudoers.lan} por el DNS del
DC}{captura-dig-web.png}{Captura de pantalla mostrando la salida de \texttt{dig web.sudoers.lan} con la respuesta \texttt{A 192.168.0.2}.}

\subsection{Registros de los equipos de administradores}

Los nombres de los miembros del equipo (\texttt{david.sudoers.lan} y
\texttt{nicolas.sudoers.lan}) tambi\'en resuelven hacia el \textbf{proxy}
(\texttt{192.168.0.2}), y no hacia las m\'aquinas de sus due\~nos: como la
puerta de entrada es \'unica, el backend puede cambiar de m\'aquina o de
puerto sin que los nombres ni los clientes cambien. Los registros se crean con
el script \textbf{idempotente} \texttt{dc/dns-records.sh}, que usa
\texttt{samba-tool dns add} con autenticaci\'on Kerberos obligatoria:

\begin{lstlisting}[style=terminal]
$ sudo kinit administrator
$ sudo bash dc/dns-records.sh
==> Servicios web alojados en m\'aquinas de admins (via proxy) ...
  ok    david.sudoers.lan    ->  192.168.0.2  (ya existe)
  ok    nicolas.sudoers.lan  ->  192.168.0.2  (ya existe)
\end{lstlisting}

\capturaAPA{fig:dns-records}{Creaci\'on idempotente de los registros de los
administradores}{captura-dns-records.png}{Captura de pantalla mostrando la salida de \texttt{sudo bash dc/dns-records.sh} con los registros \texttt{david} y \texttt{nicolas} apuntando a \texttt{192.168.0.2}.}

\subsection{Registros SRV de descubrimiento}

Los registros \textbf{SRV} son las ``gu\'ias de servicios'' del dominio
\citep{rfc2782}: le dicen a cualquier cliente d\'onde est\'a el LDAP, el
Kerberos o el DC. Sin ellas, las estaciones no encuentran el dominio. El
aprovisionamiento de Samba los gener\'o autom\'aticamente:

\begin{lstlisting}[style=terminal]
$ dig _ldap._tcp.dc._msdcs.sudoers.lan SRV
;; ANSWER SECTION:
_ldap._tcp.dc._msdcs.sudoers.lan. 900 IN SRV 0 100 389 dc1.sudoers.lan.
\end{lstlisting}

\capturaAPA{fig:dig-srv}{Registros SRV de descubrimiento del dominio generados
por el DC}{captura-dig-srv.png}{Captura de pantalla mostrando la salida de \texttt{dig \_ldap.\_tcp.dc.\_msdcs.sudoers.lan SRV} con el registro SRV apuntando al puerto 389 de \texttt{dc1.sudoers.lan}.}

\subsection{Reenviador hacia Internet}

Para que los clientes resuelvan tambi\'en nombres externos, el DC funciona con
un \textit{forwarder} hacia \textbf{el router de la red}. De esta manera una
\'unica respuesta del DNS del dominio cubre tanto los nombres internos (zona
\texttt{sudoers.lan}) como los de Internet \citep{iscBIND9}:

\begin{lstlisting}[style=terminal]
$ dig google.com @192.168.0.2 +short
142.250.78.78          # resolución externa vía el forwarder 192.168.0.1
\end{lstlisting}

\capturaAPA{fig:dig-externo}{Resoluci\'on de nombres externos a trav\'es del
\textit{forwarder}}{captura-dig-externo.png}{Captura de pantalla mostrando la salida de \texttt{dig google.com @192.168.0.2} con una direcci\'on IP p\'ublica resuelta.}

\subsection{Entrega del DNS a la red (integraci\'on con DHCP)}

Para que una estaci\'on reci\'en conectada quede lista sin intervenci\'on
manual, el servidor DHCP \textbf{Kea} (en contenedor, con \texttt{network\_mode:
host}) entrega junto con la IP las opciones de red: \textbf{DNS
\texttt{192.168.0.2}}, dominio \texttt{sudoers.lan}, gateway
\texttt{192.168.0.1} y NTP \texttt{192.168.0.2}. De este modo toda m\'aquina
de la red resuelve los nombres internos contra el DNS del DC de forma
transparente.

% ============================================================================
\section{Servidor WEB: Nginx como proxy inverso}

\subsection{Despliegue como contenedor}

El servidor web se despleg\'o como \textbf{contenedor Docker} del proyecto
(servicio \texttt{web} del \texttt{docker-compose.yml}), a diferencia del DC
que corre nativo. La imagen es \texttt{nginx:1.27-alpine} y publica los
puertos 80 y 443 del host; el portal est\'atico y los certificados TLS llegan
por vol\'umenes:

\begin{lstlisting}[style=terminal]
web:
  build: ./services/web          # nginx:1.27-alpine
  container_name: tso-web
  ports:
    - "80:80"                    # redirige a HTTPS
    - "443:443"                  # HTTPS (termina TLS)
  volumes:
    - ./services/web/public:/usr/share/nginx/html:ro
    - ./services/web/tls:/etc/nginx/tls:ro      # CA + cert wildcard
  restart: always                # la puerta web debe estar siempre arriba
\end{lstlisting}

\capturaAPA{fig:compose-web}{Servicio \texttt{web} definido en el
\texttt{docker-compose.yml}}{captura-compose-web.png}{Captura de pantalla mostrando el bloque \texttt{web:} del \texttt{docker-compose.yml} con los puertos 80/443, los vol\'umenes del portal y de TLS, y la directiva \texttt{restart: always}.}

\subsection{El portal y los virtual hosts}

Nginx enruta por \textbf{nombre} cada petici\'on: el portal de la intranet
responde a \texttt{sudoers.lan}, \texttt{web.sudoers.lan} y
\texttt{www.sudoers.lan} desde el directorio est\'atico, mientras el resto de
los \textit{virtual hosts} (\texttt{print}, \texttt{portainer}, los miembros)
son \textit{proxies} hacia sus backends. Adem\'as, \textbf{todo HTTP (80)
redirige a HTTPS (443)}: la web nunca se sirve en claro.

\begin{lstlisting}[style=terminal]
server {
    listen 80;                                 # HTTP: solo redirige
    server_name sudoers.lan web.sudoers.lan www.sudoers.lan;
    return 301 https://$host$request_uri;
}

server {
    listen 443 ssl;                            # HTTPS: portal
    http2 on;
    server_name sudoers.lan web.sudoers.lan www.sudoers.lan;
    ssl_certificate     /etc/nginx/tls/server.crt;
    ssl_certificate_key /etc/nginx/tls/server.key;
    root  /usr/share/nginx/html;
    index index.html;
}
\end{lstlisting}

\capturaAPA{fig:nginx-conf}{Virtual host del portal y redirecci\'on HTTP a
HTTPS en \texttt{services/web/nginx.conf}}{captura-nginx-conf.png}{Captura de pantalla mostrando el contenido de \texttt{services/web/nginx.conf} con el bloque \texttt{server} del portal y la directiva \texttt{return 301 https://\$host\$request\_uri}.}

\subsection{HTTPS con CA interna}

El sitio se publica por \textbf{HTTPS} con un certificado \textit{wildcard}
\texttt{*.sudoers.lan} firmado por una \textbf{CA interna} (``SUDOERS.LAN
CA''), generada con el script \texttt{services/web/tls-gen.sh} usando
\texttt{openssl}:

\begin{lstlisting}[style=terminal]
$ sudo bash services/web/tls-gen.sh
=> Generando CA interna ...
=> Generando certificado para *.sudoers.lan ...
============================================
  CA:       ca.crt (instal\'a en los clientes para confiar)
  Server:   server.crt / server.key  (*.sudoers.lan, 825 d\'ias)
============================================
$ docker compose up -d --build web
\end{lstlisting}

\capturaAPA{fig:tls-gen}{Generaci\'on de la CA interna y del certificado
\textit{wildcard}}{captura-tls-gen.png}{Captura de pantalla mostrando la salida de \texttt{sudo bash services/web/tls-gen.sh} con la creaci\'on de la CA y del certificado \texttt{*.sudoers.lan}.}

El certificado cubre los SAN \texttt{sudoers.lan}, \texttt{*.sudoers.lan},
\texttt{web}, \texttt{www}, \texttt{print}, \texttt{portainer}, \texttt{mail},
\texttt{webmail}, \texttt{dashboard} y los subdominios de los miembros. Como
la CA es interna, cada cliente debe instalar \texttt{ca.crt} como CA de
confianza del sistema para evitar la advertencia del navegador. El acceso por
HTTPS se verifica localmente con \texttt{curl}:

\begin{lstlisting}[style=terminal]
$ curl -skI https://web.sudoers.lan
HTTP/1.1 200 OK
Server: nginx
...
\end{lstlisting}

\capturaAPA{fig:curl-web}{Acceso HTTPS al portal servido por Nginx}{captura-curl-web.png}{Captura de pantalla mostrando la salida de \texttt{curl -skI https://web.sudoers.lan} con el c\'odigo \texttt{200 OK} y el encabezado \texttt{Server: nginx}.}

\subsection{Proxy a los servicios: impresi\'on y Portainer}

Los servicios contenedores sin exposi\'on p\'ublica se alcanzan a trav\'es
del \texttt{proxy}: \texttt{print.sudoers.lan} y \texttt{portainer.sudoers.lan}
resuelven (DNS) hacia \texttt{192.168.0.2} y Nginx los reenv\'ia a los
contenedores \texttt{tso-print} y \texttt{tso-portainer} por la red interna de
Docker, sin publicar puertos adicionales en el host:

\begin{lstlisting}[style=terminal]
server {
    listen 443 ssl;
    server_name print.sudoers.lan;
    location / {
        resolver 127.0.0.11 valid=10s ipv6=off;   # DNS interno de Docker
        set $cups_upstream http://tso-print:631;
        proxy_pass $cups_upstream;
    }
}
\end{lstlisting}

Solo el proxy conoce la direcci\'on interna de cada backend; desde la LAN el
\'unico puerto visible es el 443 del proxy \citep{nginxDocs}.

\subsection{Subdominios de los administradores (david / nicolas)}

David y Nicol\'as alojan su propio servicio web \textbf{en sus m\'aquinas}
(\texttt{192.168.0.51} y \texttt{192.168.0.52}, puerto \texttt{8080}). Su DNS
apunta al proxy (\texttt{192.168.0.2}) y el proxy alcanza la m\'aquina por la
red plana:

\begin{lstlisting}[style=terminal]
server {
    listen 443 ssl;
    server_name david.sudoers.lan;
    location / {
        resolver 127.0.0.11 valid=10s ipv6=off;
        set $david_upstream http://192.168.0.51:8080;
        proxy_pass $david_upstream;
        proxy_connect_timeout 5s;
    }
}
\end{lstlisting}

\begin{lstlisting}[style=terminal]
$ dig david.sudoers.lan +short          # el cliente resuelve el proxy
192.168.0.2
$ curl -sk https://david.sudoers.lan    # y el proxy responde la m\'aquina
... (contenido del servicio de David) ...
\end{lstlisting}

\capturaAPA{fig:curl-david}{Servicio de la m\'aquina de David publicado por el
proxy}{captura-curl-david.png}{Captura de pantalla mostrando la resoluci\'on \texttt{dig david.sudoers.lan +short} con \texttt{192.168.0.2} y la respuesta de \texttt{curl -sk https://david.sudoers.lan} con el contenido del servicio del administrador.}

Para el cliente no importa el puerto interno: afuera solo existe la puerta
est\'andar \texttt{443} y los nombres quedan libres de direcciones y puertos.

% ============================================================================
\section{Pruebas de integraci\'on}

\subsection{Resoluci\'on y acceso desde un cliente de la red}

Un equipo cliente de la red \texttt{192.168.0.0/24} (con el DNS
\texttt{192.168.0.2} entregado por DHCP) consulta el dominio y accede al
portal sin configuraci\'on manual:

\begin{lstlisting}[style=terminal]
$ dig web.sudoers.lan +short             # desde el cliente
192.168.0.2
$ curl -sk https://web.sudoers.lan | head -5
<!DOCTYPE html>
<html lang="es">
... (portal de la intranet) ...
\end{lstlisting}

\capturaAPA{fig:cliente-dig}{Resoluci\'on de nombres desde un cliente de la
red}{captura-cliente-dig.png}{Captura de pantalla mostrando la salida de \texttt{dig web.sudoers.lan} ejecutada en un cliente, respondiendo \texttt{192.168.0.2}.}

\capturaAPA{fig:cliente-web}{Acceso HTTPS al portal desde el navegador del
cliente}{captura-cliente-web.png}{Captura de pantalla del navegador de un cliente mostrando el portal de la intranet cargado por \texttt{https://web.sudoers.lan} con el candado de confianza de la CA interna.}

La misma validaci\'on se repite para el portal, la consola de impresi\'on, el
panel de Portainer y las p\'aginas de los miembros, confirmando que
\textbf{DNS + proxy web} funcionan como un \'unico sistema de nombres y
publicaci\'on.

\subsection{Resumen de decisiones t\'ecnicas}

La tabla \ref{tab:decisiones} resume las decisiones m\'as relevantes de la
configuraci\'on y su justificaci\'on.

\begin{table}[H]
  \centering
  \caption{Decisiones t\'ecnicas y justificaci\'on}
  \label{tab:decisiones}
  \contenidoConFuente{%
    \begin{tablaAPA}
      \begin{tabular}{@{}p{5cm}p{8.5cm}@{}}
      \toprule
      \textbf{Decisi\'on} & \textbf{Justificaci\'on} \\
      \midrule
      DNS del dominio en el DC (SAMBA\_INTERNAL) & El AD exige que el DC sea quien sirva la zona \texttt{sudoers.lan} y sus registros SRV \citep{sambateam2024} \\
      Registros A apuntan SIEMPRE al proxy (\texttt{192.168.0.2}) & Una \'unica puerta: los backends cambian sin cambiar los nombres \citep{msft_activedirectory} \\
      Nginx en contenedor (docker-compose) & La web es un servicio m\'as del compose; solo el DC corre nativo \\
      TLS con CA interna + wildcard \texttt{*.sudoers.lan} & Cifra la intranet sin depender de una CA p\'ublica ni de dominios externos \\
      HTTP $\rightarrow$ HTTPS (301) y \texttt{ssl\_reject\_handshake} & Proh\'ibe tr\'afico en claro y el handshake a nombres desconocidos \\
      \texttt{dc/dns-records.sh} idempotente & Se puede ejecutar varias veces sin duplicar registros (RECORD\_ALREADY\_EXISTS) \\
      \textit{Forwarder} hacia \texttt{192.168.0.1} & Resuelve Internet reutilizando la salida del router \\
      \bottomrule
      \end{tabular}
    \end{tablaAPA}%
  }{Elaboraci\'on propia.}
\end{table}